\section{A completed arithmetic application: exponent six}
\label{sec:census}

This section re-exposes and validates the proof in the earlier exponent-six report. The new work here is the independent review and complete replay of its decisive finite tests.\cite{priorSix}

\begin{theorem}[Completion of the exponent-six census]
There are exactly 432 imaginary quadratic number fields whose ideal class group has exact exponent six. The largest absolute fundamental discriminant among them is $5\,761\,140$. There are exactly 1,236 imaginary quadratic orders of exact class-group exponent six, with largest absolute discriminant $43\,522\,752$.
\end{theorem}

The group exponent is the least positive integer killing every ideal class. The proof uses Dirichlet GRH as supplied by the Lean source. It also uses two published inputs: the exhaustive field census through
\[
 X_0=3.1\cdot10^{20},
\]
and the quadratic-GRH upper bound
\[
 |D|\le X_1=2.5\cdot10^{25}
\]
for exponent-six fields. Both are in Elsenhans--Kl\"uners--Nicolae (EKN). Their original ERH completeness theorem excludes exponent six; the interval $(X_0,X_1]$ requires the extra argument below. The order assertion then uses the conductor classification of Aiken--Ross--Shvydkoy--Xue.\cite{EKN,ARSX}

\subsection{Integral norms constrain every small split prime}

Suppose that $D<0$ is fundamental, $x=|D|>X_0$, and the class group is killed by six. If a rational prime $p$ splits, write $p\mathcal O_K=\mathfrak p\overline{\mathfrak p}$. Then $\mathfrak p^6$ is principal, generated by an algebraic integer of the form $(a+b\sqrt D)/2$. The coefficient $b$ is nonzero: a rational generator would have equal valuations at the two conjugate prime ideals. Taking absolute values of the coefficients yields
\begin{equation}\label{eq:norm6}
 a_p^2+x b_p^2=4p^6,\qquad a_p\ge0,\quad b_p\ge1.
\end{equation}
In particular $x\le4p^6$. Since $4\cdot2000^6<X_0$, every split prime in a putative omitted field exceeds 2000.

\subsection{A signed test with nonnegative Fourier transform}

Let $Y=5500$, $\delta=\log Y/12$, and let $g$ take the following values on successive intervals of length $\delta$, starting at zero, and vanish elsewhere:
\[
 21,14,10,6,4,1,-1,-4,-6,-10,-14,-21.
\]
Define its normalized autocorrelation and Fourier transform by
\[
 F(u)=\frac{\int_\R g(v)g(v+u)\,dv}{\int_\R g(v)^2\,dv},
 \qquad h(t)=\int_\R F(u)e^{itu}\,du.
\]
The function $F$ is even, piecewise linear, supported on $[-\log Y,\log Y]$, and satisfies $F(0)=1$, $|F|\le1$. Its transform $h$ is nonnegative on $\R$, because it is a squared Fourier modulus divided by $\|g\|_2^2$. The step-function construction also makes $h$ entire with quadratic decay on each fixed horizontal strip.

For the primitive odd quadratic character $\chi(n)=(D/n)$, the Guinand--Weil formula in this normalization is
\begin{align}
 \sum_\rho h\bigl((\rho-1/2)/i\bigr)
 ={}&\log(x/\pi)
 +\frac1{2\pi}\int_\R h(t)\Re\frac{\Gamma'}\Gamma(3/4+it/2)\,dt\notag\\
 &-2\sum_{n\ge2}\frac{\chi(n)\Lambda(n)}{\sqrt n}F(\log n).
 \label{eq:explicit6}
\end{align}
The supplied quadratic GRH makes the left side nonnegative. The analytic identity and its test-function hypotheses are standard; the cited Fourier-optimization paper states a compatible version as Lemma 9.\cite{Carneiro}

Write $C$ for the maximum absolute slope of $F$, and $\psi=\Gamma'/\Gamma$. Fourier inversion and the integral formula for $\psi$ give
\[
 \frac1{2\pi}\int_\R h(t)\Re\psi(3/4+it/2)\,dt
 =\psi(3/4)+2\int_0^\infty
   \frac{e^{-3u/2}}{1-e^{-2u}}(1-F(u))\,du
 \le\psi(3/4)+\frac{14}9 C.
\]
Indeed $0\le1-F(u)\le Cu$, and
\[
 \sum_{j\ge0}(2j+3/2)^{-2}\le\frac49+\frac13=\frac79.
\]
The elementary bound $\log\pi-\psi(3/4)>2$ follows from
$\log\pi>1.1$, $\gamma>1/2$, and
$\psi(1)-\psi(3/4)>1/3+1/14$. Therefore
\begin{equation}\label{eq:prime-test6}
 \log x>
 2\sum_{n\ge2}\frac{\chi(n)\Lambda(n)}{\sqrt n}F(\log n)
 +2-\frac{14}9 C.
\end{equation}

The signed values of $g$ are significant. Although $h$ is nonnegative, $F$ can be negative near the end of its support. An inert prime, for which $\chi(p)=-1$, can then contribute positively to the right side. This combines a positive spectral test with an arithmetic exclusion pattern.

\subsection{At least 28 split primes, but at most 18 coefficients}

Use the all-inert model
\[
 R_p=2\sum_{p^k\le Y}\frac{(-1)^k}{\sqrt{p^k}}F(k\log p),
 \quad S=\sum_{p\le Y}R_p\log p,
 \quad B=\max(0,\max_{p\le Y}R_p).
\]
Removing a ramified prime loses at most $R_p\log p$, so the total loss is at most $B\log x$. Every split prime exceeds 2000, and its square exceeds $Y$; changing it from inert to split loses
\[
 L_p=-4\frac{\log p}{\sqrt p}F(\log p).
\]
If $N_Y$ split primes occur below $Y$, and $L$ bounds all these losses from above, \eqref{eq:prime-test6} implies
\[
 N_Y L>S+2-\frac{14}9C-(1+B)\log X_1.
\]
The accompanying rational checker proves, with strict enclosures,
\begin{equation}\label{eq:exact-margins}
 \max_{2000<p\le5500}L_p<0.216,
 \qquad S+2-\frac{14}9C-(1+B)\log X_1>5.84.
\end{equation}
Because $27\cdot0.216=5.832<5.84$, one obtains $N_Y\ge28$.

On the other hand, \eqref{eq:norm6} and
\[
 4\cdot5500^6<19^2X_0
\]
give $1\le b_p\le18$. Two distinct split primes $p<q$ must therefore share the same $b$.

The checker obtains rigorous logarithm bounds from the positive series for $2\operatorname{arctanh}z$, after range reduction to $1\le u\le2$, with a rational geometric tail. Square-root bounds use integer square roots. The knot containing $\log(p^k)$ is located by comparing $(p^k)^{12}$ with powers of $Y$. No floating-point sign decision enters \eqref{eq:exact-margins}.

\subsection{A finite certificate eliminates the collision}

For two equations with the same $b$, set $A=a_p$, $A'=a_q$. Their squares differ by a multiple of four, so $A,A'$ have the same parity. Put
\[
 u=(A'-A)/2,\quad v=(A'+A)/2,\quad
 N=q^6-p^6=uv.
\]
Here $0<u\le v$, $A=v-u$, and
\begin{equation}\label{eq:candidate6}
 M=4p^6-(v-u)^2=x b^2>X_0.
\end{equation}
For a divisor $u\mid N$ with $u\le\sqrt N$, positivity of $M$ is equivalent to $u>q^3-p^3$. Hence the search is finite:
\[
 p<q\le5500,\quad 4p^6>X_0,\quad
 u\mid(q^6-p^6),\quad q^3-p^3<u\le\lfloor\sqrt{q^6-p^6}\rfloor.
\]
For any odd prime $r\le2000$, a nonzero quadratic-residue value of $-M$ modulo $r$ forces $r\nmid bD$ and $(D/r)=1$, since $-M=Db^2$. That would be a forbidden small split prime. It is therefore enough to produce such a witness for every candidate $M>X_0$.

Both independent implementations factor
\[
 q^6-p^6=(q-p)(q+p)(q^2+pq+p^2)(q^2-pq+p^2)
\]
by trial division, generate all divisors, and apply exact residue tests. Every factor is at most $3\cdot5500^2$, so the stated prime-factorization range is sufficient. Python uses arbitrary-precision integers. C++ uses 128-bit products; its floating-point square-root seed is corrected by integer comparisons before use.

\begin{center}
\begin{tabular}{@{}lrr@{}}
\toprule
Quantity & Python replay & C++ replay\\
\midrule
Eligible primes &414&414\\
Prime pairs &85,491&85,491\\
Divisors in the chosen interval &10,144,433&10,229,924\\
Candidates $M>X_0$ &9,393,302&9,393,302\\
Candidates without a witness &\textbf{0}&\textbf{0}\\
\bottomrule
\end{tabular}
\end{center}

The C++ interval includes its lower endpoint, adding exactly one divisor with $M=0$ for each prime pair. This explains the difference of 85,491 in the intermediate counts. Both checks agree on the relevant candidates. In the C++ replay, primes from 19 through 151 already suffice to eliminate them all. This proves the finite lemma and contradicts the forced collision.

\subsection{What has been completed and what is reused}

The omitted fundamental-discriminant interval is empty. Combining this with EKN's published census proves the 432-field statement. The included field list contains 432 distinct entries and has the stated maximum. The prior order-table checker verifies the author table's exact Git blob identity, 1,236 distinct order discriminants, and equality of its 432 fundamental entries with the field list. The remaining 804 orders are nonmaximal. The conductor theorem in the cited paper supplies exhaustiveness for orders.\cite{ARSX,orderData}

This is a concrete positive answer to the request for something beyond a familiar GRH consequence: the cited literature still had a finite interval problem after assuming ERH, and the integral collision argument closes it. It is not a new theorem in the upstream Lean library. The present package includes the earlier checkers unchanged, explicit attribution, fresh logs, and a separate statement of the published inputs. We did not recompute EKN's entire initial census or claim an end-to-end Lean proof of this classification.
